\exercisesfor{4}

In the following exercises the letter G denotes a group.

\begin{enumerate}
\def\labelenumi{\arabic{enumi}.}
\item
  Let \((\mathbf{G}_n)\) be an integral central filtration on \(G\) . Show that the Lie algebra \(\mathrm{gr}(G)\) is generated by \(\mathrm{gr}_1(G)\) if and only if \(\mathbf{G}_n = \mathbf{G}_{n+1}.\mathbf{C}^n\mathbf{G}\) for all \(n \geqslant 1\) . In that case, show that \(\mathbf{G}_n = \mathbf{G}_m.\mathbf{C}^n\mathbf{G}\) for \(m > n\) and deduce that \((\mathbf{G}_n) = (\mathbf{C}^n\mathbf{G})\) if there exists an integer \(m\) such that \(\mathbf{G}_m = \{\varepsilon\}\) .
\item
  Let \((\mathbf{G}_{\alpha})\) be a real filtration on a group \(\mathbf{G}\) and let \(v\) be the corresponding order function. Let \(\mathbf{H}\) be a subgroup of \(\mathbf{G}\) .
\end{enumerate}

\begin{enumerate}
\def\labelenumi{(\alph{enumi})}
\item
  Let \(H_{\alpha} = H \cap G_{\alpha}\) . Show that \((\mathrm{H}_{\alpha})\) is a real filtration of H and that the corresponding order function is the restriction of v to H. \(H_{\alpha}^{+} = H \cap G_{\alpha}^{+}\) and \(\mathrm{gr}(\mathrm{H})\) is identified with a graded subgroup of \(\mathrm{gr}(\mathrm{G})\) .
\item
  Suppose that H is normal and that \(v(G) \cap \mathbf{R}\) is a discrete subset of R. We write \((G/H)_{\alpha} = (G_{\alpha}H)/H\) . Show that \(((G/H)_{\alpha})\) is a real filtration of G/H and that the corresponding order function \(v_{G/H}\) is given by the formula
\end{enumerate}

\[
v _ {\mathrm{G/H}} (x) = \operatorname * {S u p} _ {y \in x} v (y).
\]

Show that, for all \(\alpha \in \mathbf{R}\) , there is an exact sequence

\[
0 \rightarrow \mathrm{gr} _ {\alpha} (\mathrm{H}) \rightarrow \mathrm{gr} _ {\alpha} (\mathrm{G}) \rightarrow \mathrm{gr} _ {\alpha} (\mathrm{G} / \mathrm{H}) \rightarrow 0.
\]

\begin{enumerate}
\def\labelenumi{(\alph{enumi})}
\setcounter{enumi}{2}
\tightlist
\item
  With the hypotheses of \((b)\) , suppose that \((G_{\alpha})\) is a central filtration. Show that so are the filtrations on H and G/H induced by \((G_{\alpha})\) and that the Lie algebra \(\mathrm{gr}(\mathrm{G}/\mathrm{H})\) is identified with the quotient of \(\mathrm{gr}(\mathrm{G})\) by the ideal \(\mathrm{gr}(\mathrm{H})\) .
\end{enumerate}

¶ 3. Let \((\mathrm{H}_{\alpha})\) be a real filtration on a group H and let \(v_{\mathrm{H}}\) be the corresponding order function. Suppose that G operates on H on the left and that, for all \(g \in G\) , the mapping \(h \mapsto g(h)\) is an automorphism of the filtered group H.

\begin{enumerate}
\def\labelenumi{(\alph{enumi})}
\tightlist
\item
  If \(\alpha \in \mathbf{R}\) , we denote by \(G_{\alpha}\) the set of \(g \in G\) such that
\end{enumerate}

\[
v _ {\mathrm{H}} (h ^ {- 1} g (h)) \geqslant v _ {\mathrm{H}} (h) + \alpha \quad \text { for   all } h \in \mathrm{H}.
\]

Show that \((\mathbf{G}_{\alpha})\) is a real filtration of \(\mathbf{G}\) and that \(\mathbf{G}_{\alpha} = \mathbf{G}\) if \(\alpha \leqslant 0\) . Let \(v\) denote the corresponding order function.

\begin{enumerate}
\def\labelenumi{(\alph{enumi})}
\setcounter{enumi}{1}
\tightlist
\item
  Suppose that the filtration \((\mathrm{H}_{\alpha})\) is central and that \(\mathrm{G}_0^+ = \mathrm{G}\) . Show that \((\mathrm{G}_{\alpha})\) is a central filtration of G. If \(\xi \in \mathrm{gr}_{\alpha}(\mathrm{G})\) and \(\eta \in \mathrm{gr}_{\beta}(\mathrm{H})\) , let \(g\) (resp. \(h\) ) be a representative of \(\xi\) (resp. \(\eta\) ) in \(\mathrm{G}_{\alpha}\) (resp. \(\mathrm{H}_{\beta}\) ); show that the image of \(h^{-1}g(h)\) in \(\mathrm{gr}_{\alpha + \beta}(\mathrm{H})\) is independent of the choice of \(g\) and \(h\) ; if it is denoted by \(\mathrm{D}_{\xi}(\eta)\) , show that \(\mathrm{D}_{\xi}\) can be extended to a derivation of the Lie algebra \(\mathrm{gr}(\mathrm{H})\) of degree \(\alpha\) and that \(\xi \mapsto \mathrm{D}_{\xi}\) is a homomorphism of the Lie algebra \(\mathrm{gr}(\mathrm{G})\) into the Lie algebra of derivations of \(\mathrm{gr}(\mathrm{H})\) . If \(v_{\mathrm{H}}(\mathrm{H}) \cap \mathbf{R}\) is contained in a discrete subgroup \(\Gamma\) of \(\mathbf{R}\) , so is \(v(\mathrm{G}) \cap \mathbf{R}\) and, for all \(\alpha \in \mathbf{R}\) , the mapping \(\xi \mapsto \mathrm{D}_{\xi}\) defined above is injective.
\end{enumerate}

\begin{enumerate}
\def\labelenumi{\arabic{enumi}.}
\setcounter{enumi}{3}
\item
  Let H be a nilpotent group of class c and let G be the group of automorphisms of H which act trivially on \(\mathrm{H}/(\mathrm{H},\mathrm{H})\) . Show that G is nilpotent of class \(\leqslant c - 1\) . (Apply Exercise 3 to the lower central series of H and note that G operates trivially on \(\mathrm{gr}(\mathrm{H})\) .) Show that, if H is a finite p-group (p prime), so is G (same method).
\item
  Let \(\mathbf{K}\) be a commutative field and let \(\mathbf{L}\) be a finite Galois extension of \(\mathbf{K}\) with Galois group \(\mathbf{G}\) . Let \(v_{\mathbf{L}}\) be a valuation of \(\mathbf{L}\) with values in \(\mathbf{Z}\) (Commutative Algebra, Chapter VI, § 3, no. 2) which is invariant under \(\mathbf{G}\) . If \(g \in \mathbf{G}\) , we write
\end{enumerate}

\[
v (g) = \sup _ {x \in L ^ {\bullet}} v _ {L} \left(\frac {g (x) - x}{x}\right).
\]

Show that v is the order function of a separated integral filtration \((G_{n})\) on G such that \(G_{0} = G\) and that the restriction of this filtration to \(G_{1}\) is central (apply Exercise 3, taking H to be the ideal of \(v_{L}\) ). Show that \(G_{1}\) is a p-group (resp. reduces to the identity element) if the residue field of L is of characteristic p \textgreater{} 0 (resp. of characteristic zero); when the fields L and K have the same residue field, \(G_{1}\) is the kernel of the homomorphism \(\phi\) defined in Commutative Algebra, Chapter VI, § 8, Exercise 11 (b).

¶ 6. Let K{[}G{]} be the algebra of G over K and let I be the kernel of the canonical homomorphism K{[}G{]} → K (``augmentation ideal''). Then


\(K[G] = K \oplus I\) and \(I\) admits as basis the family of \(g - 1\) where \(g \in G - \{e\}\) .

\begin{enumerate}
\def\labelenumi{(\alph{enumi})}
\item
  If \(n\) is an integer \(\geqslant 0\) , we denote by \(I^n\) the ideal of \(K[G]\) the \(n\) -th power of \(I\) . Let \(G_n\) be the set of \(g \in G\) such that \(g - 1 \in I^n\) . Show that \((G_n)\) is an integral central filtration on \(G\) . In particular, \(G_n \supset C^n G\) for all \(n\) .
\item
  Show that if \(\mathbf{K} = \mathbf{Z}\) the mapping \(g \mapsto g - 1\) defines on taking quotients an isomorphism of \(\mathrm{G} / (\mathrm{G}, \mathrm{G})\) onto \(\mathrm{I} / \mathrm{I}^2\) . Deduce that \(\mathrm{G}_2 = \mathrm{C}^2\mathrm{G}.\dagger\)
\item
  Suppose that K is a field of characteristic zero and that G is finite. Show that \(I^{n} = I\) for all \(n \geqslant 1\) .
\item
  Suppose that K is a field of characteristic p \textgreater{} 0 and that G is a p-group. Show that \(I^{n}\{0\}\) for n sufficiently large. (Show first, using Algebra, Chapter I, § 6, no. 5, Proposition 11, that every simple K{[}G{]}-module is isomorphic to K; deduce that I is the radical of K{[}G{]} and hence is nilpotent since K{[}G{]} is of finite rank over K.)
\item
  Suppose that K = Z and that G has the following property: for all \(g \in G\) such that \(g \neq 1\) , there exists a prime number p, a p-group P and a homomorphism \(f: G \to P\) such that \(f(g) \neq e\) . Show that in that case \(\bigcap_{n} I^{n} = \{0\}\) . (Reduce it immediately to the case where G = P. By applying (d) to the field \(F_{p}\) , we see that there exists m such that \(I^{m} \subset p\) . \(Z[G]\) and, as I is a direct factor in Z{[}G{]}, this implies that \(I^{m} \subset pI\) , whence \(\bigcap_{n} I^{mn} \subset \bigcap_{n} p^{n}I\) , which reduces to 0 since I is a finitely generated Abelian group.)
\end{enumerate}

\begin{enumerate}
\def\labelenumi{\arabic{enumi}.}
\setcounter{enumi}{6}
\item
  Let G be given the filtration \((\mathrm{C}^{\mathrm{n}}\mathrm{G})\) and suppose that \(\mathrm{gr}_{1}(\mathrm{G}) = \mathrm{G}/(\mathrm{G}, \mathrm{G})\) is cyclic. Show that \(\mathrm{gr}_{n}(\mathrm{G}) = \{0\}\) for \(n \geqslant 2\) (use Proposition 5) and deduce that \(\mathrm{C}^{\mathrm{n}}\mathrm{G} = (\mathrm{G}, \mathrm{G})\) for \(n \geqslant 2\) .
\item
  Let \(G = \mathbf{S}\mathbf{L}_{2}(\mathbf{Z})\) and let \(x = \begin{pmatrix} 1 & 1 \\ 0 & 1 \end{pmatrix}, y = \begin{pmatrix} 1 & 0 \\ 1 & 1 \end{pmatrix}, w = \begin{pmatrix} 0 & 1 \\ -1 & 0 \end{pmatrix}.\)
\end{enumerate}

\begin{enumerate}
\def\labelenumi{(\alph{enumi})}
\item
  Verify the formulae \(w^4 = 1\) , \(w = xy^{-1}x\) , \(wxw^{-1} = y^{-1}\) .
\item
  If \(g = \begin{pmatrix} a & b \\ c & d \end{pmatrix}\) is an element of \(G\) , we write \(l(g) = |a| + |b|\) . Show that \(l(g) = 1\) if and only if \(g\) is of the form \(y^n w^a\) , where \(n \in \mathbf{Z}\) and \(0 \leqslant a \leqslant 3\) . If \(l(g) \geqslant 2\) , show that there exists a power \(h\) of \(x\) or \(y\) such that \(l(gh) < l(g)\) . Deduce that \(G\) is generated by \(\{x, y\}\) .
\item
  Using (a) and (b), show that G/(G, G) is generated by the image \(\xi\) of x and that \(\xi^{12} = e.\ddagger\) Deduce that \(C^{n}G = (G, G)\) for \(n \geqslant 2\) . (apply Exercise 7).
\end{enumerate}

\begin{enumerate}
\def\labelenumi{\arabic{enumi}.}
\setcounter{enumi}{8}
\tightlist
\item
  Let G be given the filtration \((C^{n}G)\).
\end{enumerate}

\begin{enumerate}
\def\labelenumi{(\alph{enumi})}
\item
  Show that the ideal of \(\operatorname{gr}(\mathbf{G})\) generated by \(\operatorname{gr}_2(\mathbf{G})\) is \(\sum_{q\geqslant 2}\operatorname{gr}_q(\mathbf{G})\) .
\item
  Let \((x_{i})_{i\in I}\) be a generating family of \(G\) and let \(m\geqslant 1\) . Suppose that, for all \((i,j)\in \mathrm{I}^2\) , \((x_i,x_j)^m\in \mathrm{C}^3\mathrm{G}\) . Show that, for all \(q\geqslant 2\) and all \(u\in \mathrm{C}^q\mathrm{G}\) , \(u^{m}\in \mathrm{C}^{q + 1}\mathrm{G}\) (use \((a)\) ).
\end{enumerate}

\begin{enumerate}
\def\labelenumi{\arabic{enumi}.}
\setcounter{enumi}{9}
\tightlist
\item
  Let \(x, y\) be in \(G\) and let \(r, s\) be two integers \(\geqslant 1\) . Suppose that \((x^r, y^s) = e\) . (a) Show that, for all \(n \geqslant 1\) , \((x, y)^{(rs)n} \in C^{n+2}G\) . (It can be assumed that \(G\) is generated by \(\{x, y\}\) ; then apply Exercise 9 (b), noting that \((x^r, y^s) \equiv (x, y)^{rs} \mod C^3G\) .)
\end{enumerate}

\begin{enumerate}
\def\labelenumi{(\alph{enumi})}
\setcounter{enumi}{1}
\tightlist
\item
  Suppose that \(x^r = y^s = e\) ; let \(t\) denote the g.c.d. of \(r\) and \(s\) . Show that \((x, y)^t \in \mathbf{C}^3\mathbf{G}\) and deduce that \((x, y)^{t^n} \in \mathbf{C}^{n+2}\mathbf{G}\) for all \(n \geqslant 1\) (same method).
\end{enumerate}

\begin{enumerate}
\def\labelenumi{\arabic{enumi}.}
\setcounter{enumi}{10}
\tightlist
\item
  Let H be a subgroup of G and let m be an integer \(\geqslant 1\) . Suppose that G is generated by a family \((x_{i})_{i\in I}\) such that \(x_{i}^{m}\in H\) for all i.
\end{enumerate}

\begin{enumerate}
\def\labelenumi{(\alph{enumi})}
\tightlist
\item
  Let H be given the filtration induced by the filtration (C \(^{n}\) G) on G and let gr(H) be identified with a graded Lie subalgebra of gr(G), cf.~Exercise 2. Show that, for all \(n \geqslant 0\) ,
\end{enumerate}

\[
m ^ {n}. \mathrm{gr} _ {n} (\mathrm{G}) \subset \mathrm{gr} _ {n} (\mathrm{H}).
\]

Deduce that, for all \(z \in \mathbf{H}.\mathbf{C}^n\mathbf{G}\) , \(z^m \in \mathbf{H}.\mathbf{C}^{n+1}\mathbf{G}\) .

\begin{enumerate}
\def\labelenumi{(\alph{enumi})}
\setcounter{enumi}{1}
\tightlist
\item
  If G is nilpotent, show that there exists an integer \(N \geqslant 0\) (depending only on the nilpotency class of G) such that \(z^{mN} \in H\) for all \(z \in G\) .
\end{enumerate}

\begin{enumerate}
\def\labelenumi{\arabic{enumi}.}
\setcounter{enumi}{11}
\item
  Suppose that G is nilpotent. Let H be a subgroup of G and let m be an integer \(\geqslant 1\) and let x, y be elements of G such that \(x^{m} \in H\) and \(y^{m} \in H\) . Show that there exists an integer \(N \geqslant 0\) such that \((xy)^{nN} \in H\) . (Apply Exercise 11 to the group generated by \(\{x, y\}\) and its intersection with H.)
\item
  \begin{enumerate}
  \def\labelenumii{(\alph{enumii})}
  \tightlist
  \item
    Let \(\mathbf{F}\) be a free group with basic family \(\{\overline{x},\overline{y}\}\) of two elements, let \(c\) be an integer \(\geqslant 2\) and let \(m\) be an integer \(\geqslant 1\) . We write \(\mathrm{F}^c = \mathrm{F} / C^c\mathrm{F}\) and denote by \(x,y\) the images of \(\overline{x},\overline{y}\) in \(\mathrm{F}^c\) . Let \(\mathrm{F}_m^c\) be the subgroup of \(\mathrm{F}^c\) generated by \(\{x^{m},y\}\) . Show that there exists an integer \(\mathrm{N}\geqslant 0\) such that \(z^{mN}\in \mathrm{F}_{m}^{c}\) for all \(z\in \mathrm{F}^c\) (use Exercise 11).
  \end{enumerate}
\end{enumerate}

\begin{enumerate}
\def\labelenumi{(\alph{enumi})}
\setcounter{enumi}{1}
\item
  Let \(\mathrm{I}^c\) (resp. \(\mathbf{I}_m^c\) ) be the normal subgroup of \(\mathrm{F}^c\) (resp. \(\mathbf{F}_m^c\) ) generated by \(y\) . Show that, if \(N\) is chosen as above and \(z \in \mathrm{I}^c\) , then \(z^{n^k} \in \mathbf{I}_m^c\) (note that \(\mathrm{F}^c / \mathrm{I}^c\) is an infinite cyclic group with generator the image of \(x\) and deduce that \(\mathbf{F}_m^c / \mathbf{I}_m^c \to \mathbf{F}^c / \mathrm{I}^c\) is injective and hence that \(\mathbf{I}_m^c = \mathbf{F}_m^c \cap \mathbf{I}^c\) ). In particular \(xy^m x^{-1} \in \mathbf{I}_m^c\) .
\item
  Suppose that G is nilpotent. Let H be a subgroup of G and let L be a normal subgroup of H. Let \(g \in G\) and let \(m \geqslant 1\) be such that \(g^{m} \in H\) . Show that, if N is sufficiently large, \(gl^{m^{N}} g^{-1} \in L\) for all \(l \in L\) . (If G is of class \textless c, choose N as in (a) and use the homomorphism \(f: F^{c} \to G\) such that \(f(x) = g, f(y) = l\) ; note that \(f(F_{m}^{c}) \subset H, f(I_{m}^{c}) \subset L\) and apply (b) above.)
\end{enumerate}

¶ 14. Let P be a set of prime numbers. An integer n is called a P-integer if it is ≠0 and all its prime factors belong to P. An element \(x \in G\) is called a P-torsion element if there exists a P-integer n such that \(x^{n} = \epsilon\) ; G is called a P-torsion (resp. P-torsion free) group if every element of G is a P-torsion element (resp. if no element of G other than e is a P-torsion element). G is called P-divisible if, for all \(x \in G\) and every P-integer n, there exists \(y \in G\) such that \(x = y^{n}\) .

Suppose that G is nilpotent.

\begin{enumerate}
\def\labelenumi{(\alph{enumi})}
\item
  Let H be a subgroup of G and let \(H_{P}\) be the set of \(x \in G\) such that there exists a P-integer n for which \(x^{n} \in H\) . Show that \(H_{P}\) is a subgroup of G (use Exercise 12) and that \((\mathrm{H}_{\mathrm{P}})_{\mathrm{P}} = \mathrm{H}_{\mathrm{P}}\) . The group \(H_{P}\) is called the P-saturation of H in G. If G is P-divisible, so is \(H_{P}\) .
\item
  Let \(\mathbf{L}\) be a normal subgroup of \(\mathbf{H}\) . Show that \(\mathbf{L}_{\mathbb{P}}\) is normal in \(\mathbf{H}_{\mathbb{P}}\) . (Use Exercise 13 (c) to prove that, if \(g \in \mathbf{H}_{\mathbb{P}}\) and \(l \in \mathbf{L}\) , then \(glg^{-1} \in \mathbf{L}_{\mathbb{P}}\) .)
\end{enumerate}

In particular, if L is normal in G, so is \(L_{P}\) and \(G/L_{P}\) is P-torsion-free. Deduce that the set of P-torsion elements of G is the smallest normal subgroup N of G such that G/N is P-torsion-free.

\begin{enumerate}
\def\labelenumi{(\alph{enumi})}
\setcounter{enumi}{2}
\item
  Suppose that G is P-torsion-free. Let \(n\) be a P-integer and let \(x, y \in G\) be such that \(x^n = y^n\) . Show that \(x = y\) . (Apply Exercise 10 (a) with \(r = s = n\) . Deduce that there exists a P-integer N such that \((x, y)^{\mathbb{N}} = e\) , whence \((x, y) = e\) since G is P-torsion-free; then \((x^{-1}y)^n = e\) , whence finally \(x = y\) .)
\item
  Let H be a subgroup of G, let L be a nilpotent group and let \(f: H \to L\) be a homomorphism. Let \(\Gamma\) be the graph of f in \(H \times L\) and let \(\Gamma_{P}\) be its P-saturation in \(G \times L\) . Show that \(\Gamma_{P}\) is contained in \(H_{P} \times L\) and that \(pr_{1}: \Gamma_{P} \to H_{P}\) is surjective if L is P-divisible and injective if L is P-torsion-free.
\end{enumerate}

Suppose that L is P-divisible and P-torsion-free. Show that f can be extended uniquely to a homomorphism \(f_{P}: H_{P} \to L\) and that the graph of \(f_{P}\) is \(\Gamma_{P}\) .

¶ 15. Preserving the notation of the above exercise, suppose that G is nilpotent. Let \(i: \mathbf{G} \to \overline{\mathbf{G}}\) be a homomorphism of G into a nilpotent group \(\overline{\mathbf{G}}\) . \((i, \overline{\mathbf{G}})\) is called a P-envelope of G if the following conditions hold:

\begin{enumerate}
\def\labelenumi{(\roman{enumi})}
\item
  \(\overline{\mathbf{G}}\) is P-divisible and P-torsion-free.
\item
  The kernel of \(i\) is the set of P-torsion elements of G.
\item
  The P-saturation of \(i(\mathbf{G})\) in \(\overline{\mathbf{G}}\) is equal to \(\overline{\mathbf{G}}\) .
\end{enumerate}

\begin{enumerate}
\def\labelenumi{(\alph{enumi})}
\item
  Let \((i, \overline{\mathbf{G}})\) be a P-envelope of G and let L be a P-divisible P-torsion-free nilpotent group. Show that, for every homomorphism \(f\colon \mathbf{G} \to \mathbf{L}\) , there exists one and only one homomorphism \(\bar{f}\colon \overline{\mathbf{G}} \to \mathbf{L}\) such that \(\bar{f} \circ i = f\) (reduce it to the case where G is P-torsion-free and use Exercise 14 (d)). Deduce that, if G has a P-envelope, \(\dagger\) this is unique up to isomorphism.
\item
  Let \((i, \overline{\mathbf{G}})\) be a P-envelope of G, let H be a subgroup of G, let \(\overline{\mathbf{H}}\) be the P-envelope of \(i(\mathbf{H})\) in \(\overline{\mathbf{G}}\) and let \(i_{\mathbf{H}}: \mathbf{H} \to \overline{\mathbf{H}}\) the homomorphism induced by \(i\) . Show that \((i_{\mathbf{H}}, \overline{\mathbf{H}})\) is a P-envelope of H.
\end{enumerate}

Suppose that H is normal in G. Then \(\bar{H}\) is normal in \(\bar{G}\) , cf.~Exercise 14 (b); if \(i_{G/H}: G/H \to \overline{G}/\overline{H}\) is the homomorphism induced by i, show that \((i_{G/H}, \overline{G}/\overline{H})\) is a P-envelope of G/H.

\begin{enumerate}
\def\labelenumi{\arabic{enumi}.}
\setcounter{enumi}{15}
\tightlist
\item
  We preserve the notation of the two previous exercises.
\end{enumerate}

\begin{enumerate}
\def\labelenumi{(\alph{enumi})}
\item
  Let \(\mathbf{S}_{\mathbf{P}}\) be the set of P-integers and let \(\mathbf{\dot{Z}}_{\mathbf{P}} = \mathbf{S}_{\mathbf{P}}^{-1}\mathbf{Z}\) be the ring of fractions of \(\mathbf{Z}\) defined by \(\mathbf{S}_{\mathbf{P}}\) (Commutative Algebra, Chapter II, § 2, no. 1). Suppose that G is nilpotent, P-torsion-free and P-divisible and let \(t \in \mathbf{Z}_{\mathbf{P}}\) , \(g \in G\) . Show that there exists one and only one element \(h\) of G such that \(h^s = g^{st}\) for all \(s \in \mathbf{S}_{\mathbf{P}}\) such that \(st \in \mathbf{Z}\) . The element \(h\) is called the \(t\) -th power of \(g\) and denoted by \(g^t\) . The mapping \(t \mapsto g^t\) is a homomorphism of \(\mathbf{Z}_{\mathbf{P}}\) into G. If \(t\) is invertible in \(\mathbf{Z}_{\mathbf{P}}\) , \(g \mapsto g^t\) is bijective. (Use Exercise 14 (c)).
\item
  Let \(\mathbf{A}\) be a commutative group and let \(i\) be the canonical mapping of \(\mathbf{A}\) into \(\mathbf{A}_{\mathbb{P}} = \mathbf{Z}_{\mathbb{P}} \times \mathbf{A}\) . Show that \((i, \mathbf{A}_{\mathbb{P}})\) is a P-envelope of \(\mathbf{A}\) .
\item
  Let \(\mathbf{G}\) (resp. \(\overline{\mathbf{G}}\) ) be the lower strict triangular group of order \(n\) over a ring \(k\) (resp. over the ring \(k_{\mathrm{P}} = k\otimes \mathbf{Z}_{\mathrm{P}}\) ) and let \(i\) be the homomorphism of \(\mathbf{G}\) into \(\overline{\mathbf{G}}\) defined by the canonical homomorphism of \(k\) into \(k_{\mathrm{P}}\) (Algebra, Chapter II, § 10, no. 4). Show that \((i,\overline{\mathbf{G}})\) is a P-envelope of \(\mathbf{G}\) .
\item
  Let G be a finite nilpotent group and hence the direct product of its Sylow \(p\) -groups \(G_p\) (Algebra, Chapter I, § 6, no. 7, Theorem 4). Let \(\overline{G}\) be the product of the \(G_p\) for \(p \notin P\) and let \(i\) be the canonical projection of G onto \(\overline{G}\) . Show that \((i, \overline{G})\) is a P-envelope of G.
\end{enumerate}

¶ 17. Suppose that G is nilpotent. Let P be a set of prime numbers and let \(\mathfrak{V}_{P}(G)\) be the set of normal subgroups of G of finite index whose index is a P-integer (cf.~Exercise 14). Let \(\mathcal{T}_{P}(G)\) be the topology defined on G by the filter base \(\mathfrak{V}_{P}(G)\) (General Topology, Chapter III, § 1, no. 2, Example).

\begin{enumerate}
\def\labelenumi{(\alph{enumi})}
\item
  Let N be a subgroup of G of finite index such that (G: N) is a P-integer and let \(N'\) be the intersection of the conjugates of N. Show that \((G: N')\) is a P-integer (use the fact that the group \(G/N'\) is the product of its Sylow groups); deduce that N is open with respect to \(\mathcal{T}_{\mathbb{P}}(G)\) .
\item
  Let H be a subgroup of G of finite index. Show that \(\mathcal{T}_{\mathrm{P}}(\mathrm{H})\) coincides with the topology induced on H by \(\mathcal{T}_{\mathrm{P}}(\mathrm{G})\) . (Same method.)
\item
  Let H be a normal subgroup of G such that G/H is isomorphic to Z and let x be a representative in G of a generator of G/H. Let \(N \in \mathfrak{V}_{\mathrm{P}}(\mathrm{H})\) and let \(N' = \bigcap_{n \in Z} x^{n} N x^{-1}\) . Show that \(N' \in \mathfrak{V}_{\mathrm{P}}(\mathrm{H})\) and that \(x N' x^{-1} = N'\) . If \(N''\) is the subgroup generated by \(N'\) and x, show that \((G: N'') = (H: N')\) . Deduce that \(\mathcal{T}_{\mathrm{P}}(\mathrm{H})\) is induced by \(\mathcal{T}_{\mathrm{P}}(\mathrm{G})\) .
\item
  Suppose that \(\mathbf{G}\) is finite. Show that, if \(\mathbf{H}\) is a subgroup of \(\mathbf{G}\) , there exists a sequence of subgroups
\end{enumerate}

\[
\mathrm{H} = \mathrm{H} _ {0} \subset \mathrm{H} _ {1} \subset \dots \subset \mathrm{H} _ {k} = \mathrm{G}
\]

such that \(\mathbf{H}_i\) is normal in \(\mathrm{H}_{i + 1}\) for \(0\leqslant i < k\) and that \(\mathrm{H}_{i + 1} / \mathrm{H}_i\) is finite or

isomorphic to \(\mathbf{Z}\) . Deduce, using \((b)\) and \((c)\) , that \(\mathcal{T}_{\mathrm{P}}(\mathrm{H})\) is induced by \(\mathcal{T}_{\mathrm{P}}(\mathrm{G})\) .

\begin{enumerate}
\def\labelenumi{(\alph{enumi})}
\setcounter{enumi}{4}
\tightlist
\item
  With the hypotheses of (d), let P' denote the complement of P in the set of prime numbers. Show that the closure of H under \(\mathcal{T}_{\mathrm{P}}(\mathrm{G})\) coincides with the P'-saturation of H in G (Exercise 14); in particular, G is Hausdorff if and only if it is P'-torsion-free.
\end{enumerate}

\begin{enumerate}
\def\labelenumi{\arabic{enumi}.}
\setcounter{enumi}{17}
\tightlist
\item
  The upper central series \((Z_{i}G)\) of the group G is defined inductively as follows:
\end{enumerate}

\begin{enumerate}
\def\labelenumi{(\roman{enumi})}
\item
  \(Z_{i}G = \{e\}\) if \(i \leqslant 0\) ;
\item
  \(Z_{i}G/Z_{i-1}G\) is the centre of \(G/Z_{i-1}G\) .
\end{enumerate}

Then \(\{e\} = Z_0G \subset Z_1G \subset \cdots\) and \(Z_1G\) is the centre of G; the \(Z_iG\) are characteristic subgroups of G.

\begin{enumerate}
\def\labelenumi{(\alph{enumi})}
\item
  Show that \(\mathbf{G}\) is nilpotent of class \(\leqslant c\) if and only if \(\mathbf{G} = \mathbf{Z}_c\mathbf{G}\) .
\item
  Show that \((\mathbf{C}^n\mathbf{G},\bar{\mathbf{Z}}_m\mathbf{G})\subset \mathbf{Z}_{m - n}\mathbf{G}.\)
\item
  Suppose that G is nilpotent and P-torsion-free (where P is a set of prime numbers, cf.~Exercise 14). Show that the \(Z_{1}G\) are P-saturated. (It suffices to verify this for \(Z_{1}G\) ; if n is a P-integer and \(g \in G\) is such that \(g^{n} \in Z_{1}G\) , then \(xg^{n}x^{-1} = g^{n}\) for all \(x \in G\) , whence \(xgx^{-1} = g\) by Exercise 14 (c) and it follows that g belongs to \(Z_{1}G\) .)
\end{enumerate}

